\documentclass[11pt,a4paper]{article}
\usepackage{iftex}
\ifPDFTeX
  \usepackage[T1]{fontenc}
  \usepackage[utf8]{inputenc}
  \usepackage{lmodern}
  \input{glyphtounicode}
  \pdfgentounicode=1
\else
  \usepackage{fontspec}
  \setmainfont{lmroman10-regular}[Extension=.otf, BoldFont=lmroman10-bold,
    ItalicFont=lmroman10-italic, BoldItalicFont=lmroman10-bolditalic, Kerning=Off]
\fi
\usepackage[margin=0.9in]{geometry}
\usepackage[hidelinks]{hyperref}
\hypersetup{pdftitle={Abhinav Kaushik - Cover Letter - NiCE}, pdfauthor={Abhinav Kaushik}}

\pagestyle{empty}
\setlength{\parindent}{0pt}
\setlength{\parskip}{10pt}
\raggedright

\begin{document}

{\Large\bfseries Abhinav Kaushik}\\
Senior Full-Stack AI Engineer\\
New Delhi, India \textbar{} +91 8700571859 \textbar{} \href{mailto:aabhinav.kaushik@gmail.com}{aabhinav.kaushik@gmail.com} \textbar{} \href{https://www.linkedin.com/in/abhinav-kaushik10/}{LinkedIn}

September 30, 2026

Hiring Team\\
NiCE Labs -- AI Innovation

\textbf{Re: AI Innovation Engineer (Requisition 11563)}

Dear Hiring Team,

I am applying to the AI Innovation pillar at NiCE Labs. The loop you describe -- spot a capability, build it, measure it against real metrics, then kill it or hand it off -- is how I already work. At Viavi I prototyped 15+ AI applications from ambiguous requirements and carried 7 into production or paid client work; the other 8 were stopped, which was the right call and the reason the 7 shipped.

My fiancée lives in Berlin and we are making our home there, so a Berlin-based team is where I want to build. Cognigy interests me because conversational agents are where the gap between a promising model and a dependable product is widest, and closing that gap is the work. I have shipped LLM systems that real users depended on, not notebooks, and I have killed enough of my own prototypes to argue from evidence rather than enthusiasm.

Two things I would bring on day one. I build agent loops at the level below the framework -- message arrays, tool calling, context management, prompt caching -- and I designed the automated evaluation and observability pipelines for our Question Answering and Summary agents, with LLM evaluators and tracing that cut human review by 43\%. I am also full-stack in practice: I built 21+ automated TypeScript and Next.js mockup generators that turn a ticket into a clickable React prototype, so I can take a concept from API to a demo UI that feels like a product without waiting on anyone, and I wire tools into agents over MCP.

Outside work I am building ThinkTree.AI, a non-profit AI learning platform used by NGO teachers to support students with ADHD. It is my own zero-to-one bet, shipped to real users and revised whenever their feedback contradicts my plan. I would welcome a conversation about the role.

Sincerely,\\[4pt]
Abhinav Kaushik

\end{document}
